\section{Выводы}
Собраны примеры корпуса по туризму из двух источников (Wikivoyage и Wikipedia, по 15000 статей), изучена их разметка и мета-информация, выделен текст, посчитана статистика. Для обоих источников есть встроенный поиск, значит корпус пригоден для дальнейших работ.

Затем по дампам собран корпус из 1\,000\,000 документов четырёх источников: Wikipedia (62\,\%), Reddit (30\,\%), Travel Stack Exchange (5\,\%) и вся англ. Wikivoyage (3,5\,\%). Сырые документы занимают 13,9~ГиБ, выделенный текст --- 5,6~ГиБ (972 млн слов, медиана 457 слов на документ). Для всех четырёх источников есть поиск (встроенный и Google с ограничением \texttt{site:}).

Выделенный текст составляет около 7\,\% объёма сырого HTML (10\,\% у Wikivoyage, 5\,\% у Wikipedia): основную часть занимают разметка, атрибуты Parsoid, ссылки, сноски и навигационные шаблоны. HTML статей Wikipedia в среднем в 1,3 раза больше, но текста в них меньше, чем в статьях Wikivoyage.

Недостатки и пути их решения:
\begin{itemize}
    \item Выборка Wikivoyage первого корпуса взята по алфавиту (от \enquote{'s-Hertogenbosch} до \enquote{Kennedy}) и не представительна. В корпусе на 1 млн документов Wikivoyage взята целиком.
    \item В дерево \texttt{Category:Tourism} Wikipedia попадают статьи, слабо связанные с туризмом (например, утечка газа \enquote{2022 Aqaba toxic gas leak}, министр \enquote{Salem bin Mohammed Al Mahrouqi}, автодорога \enquote{National Route A003 (Argentina)}); нужно ограничить глубину обхода или отфильтровать категории.
    \item В миллионе Wikipedia всё ещё больше половины документов, и чем ниже оценка туристичности, тем дальше статья от туризма. Порог отбора --- оценка 3: один признак места или три туристических слова. По выборке из 3000 отобранных статей у 30\,\% нет ни одного туристического слова --- только признак места (например, небольшие населённые пункты США). На пороге есть и вовсе не туристические статьи, набравшие 3 очка за слова в другом смысле (\emph{travelled}, \emph{attraction} в физике): о певице, о датчике угла поворота, список дел Верховного суда США. Улучшить отбор можно оценкой по ссылкам из Wikivoyage (статья, на которую ссылается путеводитель, --- точно туристическая) или классификатором, обученным на Wikivoyage.
    \item Вики-разметка разбирается регулярными выражениями, а не настоящим парсером: остаются следы редких шаблонов и таблиц. Лучше было бы брать HTML-дампы Wikimedia Enterprise, но на момент сбора они не обновлялись с марта 2025 года.
    \item В выделенном тексте остаются библиографические списки разделов \enquote{Further reading}, подписи к картинкам и маршрутные таблицы Wikivoyage; их можно удалять по заголовку раздела и классу элемента.
\end{itemize}
\pagebreak
